\section{Introduction}\label{sec:intro}
% Objective: State the question, the answer (no), and the two contributions: the verification result on six notes, and a census of what is priceable at all.

The cover page of a US structured note's pricing supplement carries an estimated
value: the issuer's own figure for what the security is worth, set below the
price the investor pays for it. One of the notes examined here was sold at
\$1{,}000 and disclosed at \$938.90. The practice traces to a sample comment
letter from the SEC's Division of Corporation Finance \citeyearpar{SEC2012},
which asked issuers to ``consider prominently disclosing'' such a figure; it is
staff guidance and not a rule, and no rule has since replaced it. What the filing
omits is how the number was reached. The model belongs to the issuer, so does
the volatility surface, and the discount curve embeds an internal funding rate
that appears nowhere in the document. An investor is given a valuation and no
means of checking it.

This paper asks whether the figure can be reproduced from outside the issuer,
using only public data.

The question is put to worst-of autocallable notes, which are the most
demanding instrument in the disclosure rather than the most convenient. Payment
tracks the weakest of three correlated equity underlyings. A contingent coupon
is paid only when every leg closes at or above its barrier on an observation
date, and principal is returned in full only if the weakest leg finishes above
a lower trigger. The call deserves particular attention, because the product's
name misdescribes it. On five of the six notes priced here redemption is not
automatic: three are redeemable at the issuer's option, and two are redeemed,
in the issuer's own words, ``if and only if the output of a risk neutral
valuation model \ldots\ indicates that redeeming on such date is economically
rational for us.'' Only one note in the sample redeems mechanically when a
level is reached. Valuing any of them calls for a joint model that carries each
leg's skew, a correlation structure, and an exercise rule solved from the
issuer's side of the contract, where the holder of the option is short the
investor's payoff and stops to minimise it.

Every input is public. Implied volatility surfaces are fitted to delayed
end-of-day option quotes from Cboe under the no-arbitrage conditions of raw
SVI, converted to Dupire local volatility, and simulated jointly against a
Treasury discount curve; the issuer's call is solved by Longstaff--Schwartz
regression on 200{,}000 paths. Before any note is priced the engine is made
to return the vanilla prices it was calibrated to, measured against the quoted
bid-ask width rather than in isolation, since an error only counts as large
relative to the price uncertainty of the market it came from. Four of the five
surfaces return a median error below 0.26 of the quoted bid-ask width, well
inside the band the market itself prices in. The S\&P 500 surface does not, and
\S\ref{sec:validation} attributes that to its exceptionally narrow quoted width
rather than to a poor fit: at 0.131 volatility points its absolute error falls
inside the range the other four surfaces span, while its quoted width is 4.7
times narrower than the next tightest in the sample.

The answer is no. On all six notes the independent valuation lands below the
disclosed figure, by \$11.41 to \$27.92 per \$1{,}000 of principal across the
five notes whose barriers are covered by quoted strikes, or 1.17 to 2.97 per
cent of the disclosed value. Inverting the question makes the size of the
problem clearer: solving for the constant funding spread over Treasuries that
would set model value equal to disclosed value returns a negative spread on
every note for which the solution brackets, from $-278$ to $-123$ basis points,
against an option-adjusted spread of 69 basis points on A-rated US corporate
debt on the pricing date. Reconciliation would require the issuer to borrow
below the Treasury curve. A funding rate the filing does not disclose is
therefore not a sufficient explanation for the difference, which is the
candidate a reader would reach for first.

What follows from this is narrower than it may appear, and the limits are worth
stating before the evidence. The result is not a finding of overpricing.
\citet{Henderson2011} established that these notes are issued above model
value, and the disclosed estimated value is itself below the issue price, so
both orderings already hold; the gap measured here is between the disclosed
value and an independent valuation, not between either and the price. Nor does this paper
identify what the residual is made of, though it narrows the field.
\S\ref{sec:inputs} takes the candidates one at a time and excludes four of
them: funding and observation timing move the value the wrong way, while the
proxied leg and the term extrapolation move the right way by too little or on
too few notes. Suboptimal exercise accounts for part of three notes and is
arithmetically unavailable for two. Two measured candidates survive, the
correlation assumption and the shape of the volatility surface, and a third,
the dynamics of that surface, is one a local-volatility model cannot test at
all. None of the three is published in the filings. Which is at work, or in
what combination, is not something six notes can settle. Six notes from
one week support a bounded claim and no population statement.

Two contributions remain. The first is the verification exercise itself, run on
the product class the nearest published work set aside: \citet{Bauer2020} check
issuer-disclosed values against standard option pricing models on German retail
certificates and state plainly that ``[m]ore complex products like
autocallables are left for future research.'' This paper prices them, and
reports the funding-spread inversion as a diagnostic that converts an
unexplained difference into a statement about what the difference cannot be.
The second is a census of what any such exercise can reach. Of 1{,}086
prospectus supplements filed on Form 424B2 in a single week, 129 are final
worst-of notes with a contingent coupon and a disclosed estimated value, and 71
of those, 55 per cent, can be priced from public option quotes at all. The
remainder reference underlyings with no listed options or no usable chain. Any
outside verification of this disclosure is bounded by that fraction, and
\S\ref{sec:data} reports where the rest of the market goes.
